\section{Big picture: what TFPT actually claims}

TFPT is presented here as a boundary-polarized closed theory of admissible vacua on the oriented
double cover. The formal theorem stack now starts from a thinner operational seed which functorially
induces the one-sided boundary datum, and the manuscript also shows that within the induced
boundary class no extra primitive discrete datum is needed. On the
one-sided branch, the self-adjoint boundary operator $B_\Sigma$ determines the Calder\'on projector,
the Calder\'on projector determines the seam polarization
$\iC$, and the canonical doubling then reconstructs the doubled scaffold $\Tprim$ together with
the primitive Hodge projector $\Pprim$. From this boundary kernel the manuscript next solves one
joint discrete admissibility problem, now best read through three decoders. The trace-balanced
carrier decoder
\[
6Y^2-Y-\mathbf 1=0,
\qquad
\tr Y=0
\]
fixes the rigid $E_3\oplus E_2$ split, hypercharge, and the one-family packet $S^+$. The global
winding decoder $[u_\Sigma]=1$ fixes the rigid family surface $X_f^\circ\cong\mathbb P^1\setminus
\mu_4$, the minimal family count $N_{\mathrm{fam}}=3$, and the admissible occupancy
$\Omega_{\mathrm{adm}}=48$. On the retained branch the seed decoder
\[
u:=\phiz
\]
simultaneously bridges $(\lambda_C,\beta_{\mathrm{rad}},\Omega_b,\theta_{13})$, while
$P'(z)=0$ fixes the intrinsic transport pole and the closed budget
$\sum_{f,j}L_{f,j}^{\mathrm{diag}}+N_\Phi=41$ closes the electromagnetic fixed point $\alpha$.
Only after that discrete theorem is the selector promoted to its full form
$\Padm=\Pprim P_{\mathrm{sing}}P_{\Theta}$, and the projected modular free energy drives the
continuous electromagnetic, geometric, transport, and vacuum closure: the geometric Hodge sector
yields the gravity/scalaron branch, and the exact transport grammar emits Yukawas, hadronic
admissibility, and strong-CP closure on the same admissible branch. The resulting theorem stack
does not merely emit a closed output package; it characterizes a canonical rigid object $\Tstar$ in
the rigid TFPT theorem class. The structural falsification map $\Ffals(\Tstar)$ stays in the main
text, while all comparison conventions live only in the appendix-level $\Rcmp(\Tstar)$ side card.

\Cref{fig:tfpt-headline-claims} compresses that statement into six headline claims. The point is
that the manuscript now presents three hard decoders inside one joint discrete theorem, one
continuous master functional, and one canonical rigid object rather than competing closure stories
with a separate readout layer.

\begin{figure}[t]
\centering
\scriptsize
\begin{tikzpicture}[
    headline/.style={draw=black!65, fill=gray!10, rounded corners, text width=11.2cm, minimum height=0.9cm, align=center, font=\bfseries},
    hard/.style={draw=blue!60!black, fill=blue!6, rounded corners, text width=3.25cm, minimum height=1.75cm, align=left, inner sep=6pt, font=\scriptsize},
    closure/.style={draw=orange!75!black, fill=orange!8, rounded corners, text width=3.25cm, minimum height=1.75cm, align=left, inner sep=6pt, font=\scriptsize},
    bridge/.style={draw=green!50!black, fill=green!8, rounded corners, text width=3.25cm, minimum height=1.75cm, align=left, inner sep=6pt, font=\scriptsize},
    note/.style={font=\footnotesize\itshape, text=black!60, align=center}
]
\node[headline] (head) at (0,0) {TFPT in one glance: three decoders and six structural payoffs};

\node[hard] (a) at (-3.95,-1.95) {\textbf{Hard: Boundary polarization}\\[2pt]$\Tbound$ fixes $(\taudbl,\iC)$, the doubled scaffold, and $\Pprim$ before any carrier split.};
\node[hard] (b) at (0,-1.95) {\textbf{Hard: Carrier decoder}\\[2pt]$6Y^2-Y-\mathbf1=0$, $\tr Y=0$ force the rigid $3+2$ split and $S^+$.};
\node[closure] (c) at (3.95,-1.95) {\textbf{Hard: Winding decoder}\\[2pt]$[u_\Sigma]=1$ fixes counting: $\mathbb P^1\setminus\mu_4$, $N_{\mathrm{fam}}=3$, and $\Oadm=48$.};

\node[closure] (d) at (-3.95,-4.2) {\textbf{Closure: Geometric Hodge sector}\\[2pt]One seam-even channel yields $\Mpl$, $\vgeo$, the Einstein term, and the scalaron branch.};
\node[closure] (e) at (0,-4.2) {\textbf{Closure: Transport \& seed}\\[2pt]$P'(z)=0$ fixes $\delta_{\mathrm{ph}}$; $u=\phiz$ and the $41$ budget close flavor and $\alpha$.};
\node[bridge] (f) at (3.95,-4.2) {\textbf{Closed outputs}\\[2pt]The branch ends in $\Tstar$, predictions, null tests, and falsifiers without a bridge layer.};

\node[note] at (0,-5.55) {$Y$ generates structure, $[u_\Sigma]=1$ generates counting, and $u$ generates bridge observables.};
\end{tikzpicture}
\caption{Headline claims of the present draft. The figure is intentionally simpler than the
dependency maps that follow: its job is to put the manuscript's decoder core on one page while
keeping the closed datum, theorem, and derived-consequence flow visible.}
\label{fig:tfpt-headline-claims}
\end{figure}

\begin{tcolorbox}[colback=tfptblue!3, colframe=tfptblue, breakable,
    title={TFPT master equation and closed indices}, colbacktitle=tfptblue]
\begin{equation}
\label{eq:tfpt-master}
\begin{aligned}
S_{\mathrm{TFPT}}
=
\operatorname{Tr}f\!\left(\frac{D_{\mathrm{rel}}+A_{\Sigma}^{\mathrm{seed}}}{\chiseed}\right)
-
\operatorname{Tr}f\!\left(\frac{D_{\mathrm{ref}}}{\chiseed}\right)
+\frac{i\pi}{2}\Delta\eta_\Sigma,
\qquad
A_{\Sigma}^{\mathrm{seed}}:=\mathcal B_{\mathrm{rel}}\oplus\nabla_{\chiseed}.
\end{aligned}
\end{equation}
\begin{equation}
\label{eq:tfpt-selectors}
\begin{aligned}
\operatorname*{Stat}_{\alpha,\chiseed,e,\wspin}S_{\mathrm{TFPT}}=0,
\qquad
\Padm:=\Pi_{\ker\Delta_{\mathrm{adm}}},
\qquad
\partial_\alpha\Gamma_{U(1)}(\alpha)=0,
\\
\operatorname{Ind}(D^+_{\mathrm{fam}})=3,
\qquad
\operatorname{Ind}(\mathcal B^+_{E_2})=1,
\qquad
\operatorname{Ind}(\mathcal B^+_{E_3})=0.
\end{aligned}
\end{equation}
\end{tcolorbox}

The hard carrier theorems below do not require the full theorem stack behind
\Cref{eq:tfpt-master,eq:tfpt-selectors}. In the present paper the master equation is the
organizing closed form: the carrier theorems are hard, the later closure theorems are internal
statements on the reconstructed branch, and the empirical interface is kept outside this theorem
box. At
this overview level the seed package is written as $A_{\Sigma}^{\mathrm{seed}}$; on the later
geometric branch it is reorganized as $A_{\Sigma}^{\mathrm{geo}}$ with
$\chiseed\mapsto\chigeo=\Lambda e^\sigma$.

\begin{remark}[Selector versus dynamics]
\label{rem:selector-versus-dynamics}
The admissibility projector $\Padm$ selects the physical sector but is not itself the full
Minkowski dynamics. The actual dynamics is carried by the relative generating functional
$Z_{\mathrm{rel}}$, its Schwinger distributions, the reconstructed local net on the admissible
Hilbert space, the running effective actions $\Gamma_k$ and
$\Gamma_{\mathrm{TFPT}}^{\mathrm{ren}}$, and the resulting scattering theory.
\end{remark}

\begin{tcolorbox}[colback=tfptblue!3, colframe=tfptblue, breakable,
    title={Closed TFPT Chain}, colbacktitle=tfptblue]
\begin{equation}
\label{eq:compiler-chain}
\begin{aligned}
\Tbound
&\;\Rightarrow\;
(\taudbl,\iC,\Pprim,[u_\Sigma],\cthree)
\;\Rightarrow\;
d^\star_{\mathrm{disc}}\\
&\;\Rightarrow\;
\Padm
\\
&\;\Rightarrow\;
\mathcal M_{d^\star_{\mathrm{disc}}}
\;\Rightarrow\;
\mathcal Z_{\mathrm{cl}}=\{x^\star_{\mathrm{cl}}\}
\;\Rightarrow\;
\Tstar,\\
\Tstar
&=
\bigl(
E_3\oplus E_2,\ X,\ Y,\ a_\Sigma,\ \Gamma_{\mathrm{grav}},\\
&\qquad
\Lambda_{\mathrm{IR}},\ x_{\mathrm{cl}}^\star
\bigr).
\end{aligned}
\end{equation}
At chain level the same architecture is read through three hard decoders:
\[
\begin{aligned}
Y&\Rightarrow \text{carrier structure},\\
[u_\Sigma]=1&\Rightarrow \text{family counting},\\
u=\phiz&\Rightarrow \text{bridge observables}.
\end{aligned}
\]
On the retained branch this further compresses to
\[
P'(z)=0\Rightarrow \delta_{\mathrm{ph}},
\qquad
\sum_{f,j}L_{f,j}^{\mathrm{diag}}+N_\Phi=41\Rightarrow \alpha.
\]
From $\Tstar$ the structural falsification card $\Ffals(\Tstar)$ remains inside the main text;
the appendix-only comparison card $\Rcmp(\Tstar)$ collects scheme conventions and empirical
ledgers.
\end{tcolorbox}

\begin{tcolorbox}[colback=tfptorange!4, colframe=tfptorange!80!black, breakable,
    title={Closed Theorem Stack in the Present Version}, colbacktitle=tfptorange!85!black]
\[
\Tbound
\;\Rightarrow\;
(\taudbl,\iC,\Pprim,[u_\Sigma],\cthree)
\;\Rightarrow\;
d^\star_{\mathrm{disc}}
\;\Rightarrow\;
\Padm
\;\Rightarrow\;
\mathcal M_{d^\star_{\mathrm{disc}}}
\;\Rightarrow\;
\mathcal Z_{\mathrm{cl}}=\{x^\star_{\mathrm{cl}}\}
\;\Rightarrow\;
\Tstar.
\]
Hard-decoder form:
\[
\begin{aligned}
6Y^2-Y-\mathbf 1=0,\ \tr Y=0
&\Rightarrow
3+2,\\
[u_\Sigma]=1
&\Rightarrow
N_{\mathrm{fam}}=3,\\
u=\phiz
&\Rightarrow
(\lambda_C,\beta_{\mathrm{rad}},\Omega_b,\theta_{13}).
\end{aligned}
\]
The manuscript contains no additional primitive ontology beyond the operational seed layer. The
falsification map $\Ffals$ is a structural side card on the main-text side; the comparison map
$\Rcmp$, optional arithmetic continuations, asymptotic expansions, and horizon language are
kept only in appendices.
\end{tcolorbox}

\begin{figure}[t]
\centering
\scriptsize
\begin{tikzpicture}[
    closed/.style={draw=blue!60!black, fill=blue!6, rounded corners, text width=2.45cm, minimum height=1.05cm, align=center},
    star/.style={draw=blue!70!black, very thick, fill=blue!12, rounded corners, text width=2.45cm, minimum height=1.05cm, align=center},
    fals/.style={draw=orange!70!black, fill=orange!8, rounded corners, text width=2.55cm, minimum height=0.95cm, align=center},
    appendix/.style={draw=green!50!black, dashed, line width=0.8pt, fill=green!4, rounded corners, text width=2.55cm, minimum height=0.95cm, align=center},
    note/.style={font=\footnotesize\itshape, text=black!70, align=center}
]
\node[closed] (a) at (-6.7,0) {boundary datum\\$\Tbound$};
\node[closed] (b) at (-3.95,0) {doubling / polarization\\$\taudbl,\iC,\Tprim$};
\node[closed] (c) at (-1.2,0) {decoder core\\$Y,\ [u_\Sigma],\ u$};
\node[closed] (d) at (1.55,0) {full projector $\Padm$\\and closure stack};
\node[star]   (e) at (4.3,0) {canonical rigid object\\$\Tstar$};
\node[fals]      (f) at (1.55,-1.85) {structural falsification\\map $\Ffals(\Tstar)$};
\node[appendix]  (g) at (4.3,-1.85)  {appendix comparison\\map $\Rcmp(\Tstar)$};
\draw[->, thick] (a) -- (b);
\draw[->, thick] (b) -- (c);
\draw[->, thick] (c) -- (d);
\draw[->, thick] (d) -- (e);
\draw[->, thick, orange!70!black] (e.south) .. controls +(0,-0.6) and +(0.5,0.4) .. (f.east);
\draw[->, dashed, line width=0.8pt, green!50!black] (e) -- (g);
\node[note] at (-1.2,-2.6) {$Y$ fixes structure, $[u_\Sigma]=1$ fixes counting, and $u$ fixes bridge observables.};
\end{tikzpicture}
\caption{Closed theorem surface with two hard-separated side maps and a decoder core.}
\label{fig:status-ladder}
\end{figure}

The remainder of the main text sharpens each arrow of this chain in turn.

\section{Primitive core and hard carrier kernel}

\subsection{Primitive core and conventions}

We use natural units $c=\hbar=1$. Relative operators are formulated in Euclidean signature and compared through a declared Wick-rotation bridge when Lorentzian language is needed. Higgs fields are denoted by $\Phi$, Hamiltonians by $\mathcal H$, and the cosmological Hubble rate by $\Hhub(t)$; this avoids overloading the letter $H$. Newton's constant is always $G_N$, while $G_{\mathrm{car}}$ is reserved for the internal carrier stabilizer. The seam is denoted by $\Seam$, geometric reflection by $\taudbl$, boundary polarization by $\iC$, carrier polarization by $\epscar$, and word parity by $\varepsilon(w)$.

Here $S\to \widetilde M$ denotes the spinor bundle on the oriented double cover, and by the same symbol we denote its restriction to
\[
\Xbulk=\widetilde M\setminus\Seam.
\]

When color-center notation is needed later, we use multiplicative notation
\[
z(w)\in\{1,\omega,\omega^2\},
\qquad
\omega=e^{2\pi i/3}.
\]
Thus center neutrality means $z(w)=1$.

\[
\Tcore^{\mathrm{kin}}
:=
\bigl(
\widetilde M\to M,\ \Seam,\ \taudbl,\ \iC,\ D_{\mathrm{ref}},\ D_{\mathrm{rel}},\ \chiseed
\bigr).
\]
When a chosen almost-commutative factorization
$\mathcal H\cong L^2(\widetilde M,S)\otimes\Hint$ is used, $S$ and $\Hint$ refer to that
choice and are not part of the kinematic core datum.
Here \enquote{relative} always means comparison with a declared reference operator or
reference-subtracted action. The canonical admissibility complex is introduced immediately below,
while any Bernstein or other spectral test profile is treated only as an analytic realization
choice rather than as part of the boundary datum. The technical tables for relative objects,
layered closure/benchmark data, and the constants atlas are collected in
Appendix~\ref{app:tech} and Appendix~\ref{app:constants}; they are kept out of the early
main-text flow so that the carrier kernel appears before the technical bookkeeping apparatus.

\begin{remark}[Boundary datum versus derived closed structure]
The primitive ontology now stops one step earlier. The one-sided boundary datum $\Tbound$
contains only one-sided operator data together with the self-adjoint boundary operator
$B_\Sigma$. Its Calder\'on projector defines the polarization
\[
\iC:=2P_C-1.
\]
From that boundary polarization and the exact double one reconstructs the deck involution
$\taudbl$, the physical quotient $M=\widetilde M/\taudbl$, the seam
$\Seam=\mathrm{Fix}(\taudbl)$, the reference/relative operator split
$(D_{\mathrm{ref}},D_{\mathrm{rel}})$, the primitive seam response
$\chiseed=\lambda_1^+(|B_\Sigma|)^{-1}$, and the doubled closed datum
$\Tmin^{\mathrm{cl}}$. The seam is therefore no longer primitive ontology but the fixed locus of
the deck involution extracted from the one-sided boundary problem.
\end{remark}

